\section*{Chapter \ref{ch:g12:income-and-wellbeing} --- Income and Wellbeing}
% figdata: figdata/grade-12/15-income-and-wellbeing/out/ (lifesat, well_key, hours_ts, five); the five countries are invented.

\begin{solution}{exo:g12:income-and-wellbeing:1}
About 56,226 dollars per person and a satisfaction of 8.1.
\end{solution}

\begin{solution}{exo:g12:income-and-wellbeing:2}
Romania is about 26\,\% richer but 2.0 points more satisfied. For that income gap the fit predicts only
\[0.93 \times \log_2(34{,}829/27{,}691) \approx 0.31 \text{ points:}\]
most of the difference has other causes.
\end{solution}

\begin{solution}{exo:g12:income-and-wellbeing:3}
Two doublings: 1.86 points. Three doublings: 2.79 points.
\end{solution}

\begin{solution}{exo:g12:income-and-wellbeing:4}
$(2{,}427 - 1{,}372)/2{,}427 \approx 43.5\,\%$.
\end{solution}

\begin{solution}{exo:g12:income-and-wellbeing:5}
$(0.8 + 0.6 + 0.7)/3 = 0.7$.
\end{solution}

\begin{solution}{exo:g12:income-and-wellbeing:6}
Income is multiplied by 4, two doublings, for 1.8 points: 0.9 per doubling. $b = 0.9/\ln 2 \approx 1.298$ and $a = 6.0 - 1.298 \ln 10{,}000 \approx -5.96$.
\end{solution}

\begin{solution}{exo:g12:income-and-wellbeing:7}
$\log_2(131{,}734/11{,}188) \approx 3.56$ doublings for 1.6 points: about 0.45 per doubling, half the full fit's 0.93. Luxembourg lies well below the line, and a line through
two points depends entirely on them; the fit through all 37 is less sensitive to any one country.
\end{solution}

\begin{solution}{exo:g12:income-and-wellbeing:8}
$0.93 \times \log_2 1.1 \approx 0.93 \times 0.1375 \approx 0.13$ points.
\end{solution}

\begin{solution}{exo:g12:income-and-wellbeing:9}
Equal weights: X $0.75$, Y $0.70$, so X ranks first. One-third and two-thirds: X $0.3 + 0.4 = 0.70$, Y $0.2 + 0.533 \approx 0.73$, so Y ranks first.
\end{solution}

\begin{solution}{exo:g12:income-and-wellbeing:10}
$786/50 \approx 15.7$ hours a week less.
\end{solution}

\begin{solution}{exo:g12:income-and-wellbeing:11}
Her income falls by $30 \times 200 = 6{,}000$ dalens. If she chose the change, the free time is worth at least 6,000 dalens to her; GDP records the fall in output, not the
gain in time.
\end{solution}

\begin{solution}{exo:g12:income-and-wellbeing:12}
Satisfaction that stops rising above some income, or a country whose income grows for decades without any rise in average satisfaction. The first has been found only for
the least happy people; the second is a 1974 finding (Easterlin 1974); later work with other data found a positive relation over time, and the question is still argued over.
\end{solution}

\begin{solution}{exo:g12:income-and-wellbeing:13}
Other things that people value, such as health, security and trust, rise with income too, and may explain part of the slope; and causation may run the other way, from
satisfied, healthy people to higher earnings. People in different countries may also use the scale differently.
\end{solution}

\begin{solution}{exo:g12:income-and-wellbeing:14}
G's output per hour is $65{,}000/1{,}800 \approx 36.1$ dollars; 400 hours are worth about 14,444 dollars. Adjusted, F has about 74,444 against G's 65,000.
\end{solution}

\begin{solution}{exo:g12:income-and-wellbeing:15}
Average log income in A is $\ln\sqrt{20{,}000 \times 60{,}000} = \ln 34{,}641$, in B $\ln\sqrt{30{,}000 \times 50{,}000} = \ln 38{,}730$. The difference in satisfaction is
$1.346 \times \ln(38{,}730/34{,}641) \approx 1.346 \times 0.112 \approx 0.15$ in B's favour. Given equally to everyone, 34,641 dollars would give A its average satisfaction
and 38,730 would give B its own.
\end{solution}

\begin{solution}{pb:g12:income-and-wellbeing:1}
\textbf{1.} A 2, B 3, C 4, D 5, E 6.

\textbf{2.} $\bar x = 4$; mean satisfaction $(5.0 + 5.7 + 6.3 + 6.9 + 7.5)/5 = 6.28$.

\textbf{3.} $\sum (x - 4) s = -10.0 - 5.7 + 0 + 6.9 + 15.0 = 6.2$; $\sum (x - 4)^2 = 10$; slope 0.62 per doubling.

\textbf{4.} $s = 6.28 + 0.62 (x - 4) = 3.80 + 0.62 x$.

\textbf{5.} $x = 7$: $6.28 + 1.86 = 8.14$.

\textbf{6.} Predictions 5.04, 5.66, 6.28, 6.90, 7.52; residuals $-0.04$, $+0.04$, $+0.02$, $0.00$, $-0.02$.

\textbf{7.} $0.62 \times \log_2 1.1 \approx 0.085$ points.

\textbf{8.} $(7.0 - 6.28)/0.62 \approx 1.16$ doublings: $16{,}000 \times 2^{1.16} \approx 35{,}800$ dollars.

\textbf{9.} $2{,}270 - 1{,}660 = 610$ hours.

\textbf{10.} No: the fit compares countries that differ in many ways; it describes how satisfaction and income move together, not what raising C's income would do.

\textbf{11.} The index is $x/7$: 0.286, 0.429, 0.571, 0.714, 0.857.

\textbf{12.} 0.350, 0.400, 0.550, 0.700, 0.825.

\textbf{13.} 0.318, 0.414, 0.561, 0.707, 0.841.

\textbf{14.} All four rank the countries A, B, C, D, E from lowest to highest: here the indicators move together.

\textbf{15.} One year raises the life index by $1/40 = 0.025$ and the composite by 0.0125. The income index must rise by 0.025, so $x$ by $0.025 \times 7 = 0.175$: income
rises by $2^{0.175} - 1 \approx 12.9\,\%$.

\textbf{16.} $0.62 \times 0.175 \approx 0.11$ points.

\textbf{17.} $\tfrac13 \Delta I_y = \tfrac23 \times \tfrac1{40}$, so $\Delta I_y = 0.05$, $\Delta x = 0.35$, and income rises by $2^{0.35} - 1 \approx 27.5\,\%$.

\textbf{18.} The bounds 1,000 and 128,000 dollars, 50 and 90 years, the use of the logarithm, and the equal weights.

\textbf{19.} Satisfaction that stops rising, or falls, between D and E; or, over time, rising incomes in these countries with no rise in satisfaction.

\textbf{20.} The fitted gain is 0.62 points per doubling of income; under the equal-weight composite, one more year of life expectancy is worth the same as a 12.9\,\% rise in
income.
\end{solution}
